%!TEX program = xelatex
%!TEX root = ../all_short.tex

\chapter{Nonparametric regression: A Summary}

\section*{Fully Bayesian Regression Recap}

\begin{itemize}
    \item In fully Bayesian regression, no specific model parameters $\overline\w^*$ are identified. Instead, the predictive distribution $p(t|\x)$ is derived.
    \item Under Gaussian assumptions:
    \begin{myequation}
    p(t|\x, \X,\t, \alpha,\beta)=\mathcal{N}(t|m(\x), \sigma^2(\x))
    \end{myequation}
    with:
    \begin{myalign}
    m(\x)&=\beta\ox^{T}\S\oX^{T}\t\\
    \sigma^{2}(\x)&=\frac{1}{\beta}+\ox^{T}\S\ox
    \end{myalign}
    where $\S=(\alpha\I+\beta\oX^T\oX)^{-1} \in\Real^{(d+1)\times (d+1)}$.
\end{itemize}

\subsection*{Equivalent Kernel}

\begin{itemize}
    \item Prediction (expectation of predictive distribution):
    \begin{myequation}
    h(\x)=m(\x)=\beta\ox^{T}\S\oX^{T}\t
    \end{myequation}
    \item Since $\oX^T\t = \sum_{i=1}^n \ox_i t_i$, we can write:
    \begin{myequation}
    h(\x)=\sum_{i=1}^{n}\beta\ox^{T}\S\ox_{i}t_{i}
    \end{myequation}
    \item Prediction is a linear combination of target values $t_i$, with weights depending on $\x$ and $\x_i$.
    \item Define \alert{equivalent kernel} $\kappa(\x_1,\x_2)=\beta\ox_1^{T}\S\ox_2$.
    \item Then:
    \begin{myequation}
    h(\x)=\sum_{i=1}^{n}\kappa(\x,\x_{i})t_{i}
    \end{myequation}
    \item With basis functions $\Vphi$:
    \begin{myequation}
    \kappa(\x_1,\x_2)=\beta\Vphi(\x_1)^{T}\S\Vphi(\x_2)
    \end{myequation}
    \item $\kappa(\x,\x_i)$ measures relevance of known target $t_i$ for predicting at $\x$.
    \item Equivalent kernel tends to assign greater relevance to targets $t_i$ corresponding to items $\x_i$ near $\x$ (localization property).
    \item Instead of introducing basis functions, we can directly use a \alert{localized kernel} defined on $\Real^d\times\Real^d$.
\end{itemize}

\section*{Kernel Regression}

\begin{itemize}
    \item Predict target for $\x$ by referring to training items, especially those closer to $\x$.
    \item Controlled by predefined \alert{kernel} function $\kappa_h(\x)$ (non-negligible only near 0).
    \item Common kernel: Gaussian (RBF) kernel:
    \begin{myequation}
    g(\x)=e^{-\frac{{\left\lVert\x\right\rVert}^2}{2h^2}}
    \end{myequation}
    \item Goal: approximate conditional expectation:
    \begin{myequation}
    \me{t|\x} = \frac{\int t\; p(\x,t)dt}{\int p(\x,t)dt}
    \end{myequation}
    \item Approximate joint distribution using kernel:
    \begin{myequation}
    p(\x,t)\approx\frac{1}{n}\sum_{i=1}^n\kappa_h(\x-\x_i)\kappa_h(t-t_i)
    \end{myequation}
    \item Assuming kernel is non-negative, has mean 0 and area 1 (probability density), we get \alert{Nadaraya-Watson} estimator:
    \begin{myequation}
    h(\x)= \frac{\sum_{i=1}^n\kappa_h(\x-\x_i)t_i}{\sum_{i=1}^n\kappa_h(\x-\x_i)} = \sum_{i=1}^n w_i(\x) t_i
    \end{myequation}
    with $w_i(\x)=\frac{\kappa_h(\x-\x_i)}{\sum_{j=1}^n\kappa_h(\x-\x_j)}$.
    \item With basis functions $\Vphi$:
    \begin{myequation}
    h(\x)=\sum_{i=1}^n w_i(\Vphi(\x)) t_i,\quad w_i(\x)=\frac{\kappa_h(\Vphi(\x)-\Vphi(\x_i))}{\sum_{j=1}^n\kappa_h(\Vphi(\x)-\Vphi(\x_i))}
    \end{myequation}
\end{itemize}

\section*{Locally Weighted Regression (LOESS)}

\begin{itemize}
    \item Improves Nadaraya-Watson by using weighted version of sum of squared differences loss.
    \item For predicting $t$ at $\x$, define local loss with weights $\psi_i(\x)=\kappa_h(\x-\x_i)$:
    \begin{myequation}
    L(\x)=\sum_{i=1}^n\psi_i(\x)(\ow^T\ox_i-t_i)^2
    \end{myequation}
    \item Minimization yields:
    \begin{myequation}
    \ow^*(\x)= (\oX^T\Psi(\x)\oX)^{-1}\oX^T\Psi(\x)\t
    \end{myequation}
    where $\Psi(\x)$ is diagonal with $\Psi(\x)_{ii}=\psi_i(\x)$.
    \item Prediction:
    \begin{myequation}
    h(\x) = \ow^*(\x)^T\ox
    \end{myequation}
\end{itemize}

\section*{Local Logistic Regression}

\begin{itemize}
    \item Same approach applied to classification using weighted cross-entropy:
    \begin{myequation}
    L(\x)=\sum_{i=1}^n\kappa_h(\x-\x_i)(t_i\log p_i-(1-t_i)\log(1-p_i))
    \end{myequation}
    with $p_i=\sigma(\ow^T\ox_i)$.
\end{itemize}

\section*{Gaussian Processes}

\subsection*{Definition}

\begin{itemize}
    \item A \alert{stochastic process} $f(\x)$ is a collection of random variables $\{f(\x): \x\in\chi\}$.
    \item It is a \alert{Gaussian process} (GP) if for any finite subset $\X=(\x_1,\ldots,\x_n)$, the values $f(\x_1),\ldots,f(\x_n)$ have a joint multivariate Gaussian distribution.
    \item GP generalizes multivariate Gaussians to infinite dimensions.
    \item Specified by:
    \begin{itemize}
        \item \alert{Mean function} $m(\x)$: for any $\x$, $\mathbb{E}[f(\x)] = m(\x)$.
        \item \alert{Covariance function} (kernel) $\kappa(\x,\x')$: $\text{Cov}(f(\x), f(\x')) = \kappa(\x,\x')$.
    \end{itemize}
    \item For any finite set $\X$, $f(\X) \sim \mathcal{N}(\Vmu_{\X}, \VSigma_{\X})$ where:
    \begin{itemize}
        \item $\Vmu_{\X} = (m(\x_1),\ldots,m(\x_n))^T$
        \item $\VSigma_{\X}[i,j] = \kappa(\x_i,\x_j)$ (Gram matrix $G_{\X}$)
    \end{itemize}
    \item $\kappa$ must be a \alert{positive definite kernel}: for any distinct $\x_1,\ldots,\x_n$ and constants $c_i$ not all zero:
    \begin{myequation}
    \sum_{i=1}^n\sum_{j=1}^n c_i c_j \kappa(\mathbf{x}_i,\mathbf{x}_j) > 0
    \end{myequation}
    \item Common kernel: RBF kernel $\kappa(\x_1,\x_2)=\sigma^2 e^{-\frac{||\x_1-\x_2||^2}{2\tau^2}}$.
    \item RBF kernel assigns higher covariance to nearby points, producing smooth functions. Smoothing increases with $\tau$.
\end{itemize}

\subsection*{Posterior Distribution (Noise-Free Case)}

\begin{itemize}
    \item GP defines a prior distribution over functions: $f \sim \mathcal{GP}(m, \kappa)$.
    \item Assume we observe $t_i = f(\x_i)$ exactly (no noise) at training points $\X$.
    \item For new points $\Z$, joint distribution of $f(\X)$ and $f(\Z)$ is:
    \begin{myalign}
    \begin{pmatrix} f(\X) \\ f(\Z) \end{pmatrix} \sim \mathcal{N}\left( \begin{pmatrix} \Vmu_{\X} \\ \Vmu_{\Z} \end{pmatrix}, \begin{pmatrix} G_{\X} & G_{\Z,\X} \\ G_{\Z,\X}^T & G_{\Z} \end{pmatrix} \right)
    \end{myalign}
    where $G_{\Z,\X}[i,j] = \kappa(\z_i,\x_j)$.
    \item Using properties of conditional Gaussians, predictive distribution $p(f(\Z)|\Z,\X,f(\X))$ is Gaussian with:
    \begin{myalign}
    \Vmu_{pr} &= \Vmu_{\Z} + G_{\Z,\X} G_{\X}^{-1} (f(\X) - \Vmu_{\X}) \\
    \VSigma_{pr} &= G_{\Z} - G_{\Z,\X} G_{\X}^{-1} G_{\Z,\X}^T
    \end{myalign}
    \item For a single test point $\z$:
    \begin{myalign}
    \mu_{pr} &= m(\z) + G_{\z,\X} G_{\X}^{-1} (\t - \Vmu_{\X}) \\
    \sigma^2_{pr} &= \kappa(\z,\z) - G_{\z,\X} G_{\X}^{-1} G_{\z,\X}^T
    \end{myalign}
    \item With zero mean assumption ($m(\x)=0$):
    \begin{myequation}
    \mu_{pr} = G_{\z,\X} G_{\X}^{-1} \t
    \end{myequation}
    \item In noise-free case, predictive mean interpolates training points exactly: for $\x_i \in \X$, $\mu_{pr}(\x_i) = t_i$ and $\sigma^2_{pr}(\x_i) = 0$.
\end{itemize}

\subsection*{Gaussian Process Regression with Gaussian Noise}

\begin{itemize}
    \item Realistic assumption: $t_i = f(\x_i) + \varepsilon$ with independent Gaussian noise $\varepsilon \sim \mathcal{N}(0, \sigma_f^2)$.
    \item Then $\text{Cov}(t_i, t_j) = \kappa(\x_i,\x_j) + \sigma_f^2 \delta_{ij}$.
    \item Covariance matrix for training points becomes:
    \begin{myequation}
    \hat\VSigma_{\X} = G_{\X} + \sigma_f^2 \I
    \end{myequation}
    \item For new points $\Z$, joint distribution:
    \begin{myalign}
    \begin{pmatrix} \t \\ f(\Z) \end{pmatrix} \sim \mathcal{N}\left( \begin{pmatrix} \Vmu_{\X} \\ \Vmu_{\Z} \end{pmatrix}, \begin{pmatrix} \hat\VSigma_{\X} & G_{\Z,\X} \\ G_{\Z,\X}^T & G_{\Z} \end{pmatrix} \right)
    \end{myalign}
    \item Predictive distribution $p(f(\Z)|\Z,\X,\t)$ is Gaussian with:
    \begin{myalign}
    \hat\Vmu_{pr} &= \Vmu_{\Z} + G_{\Z,\X} \hat\VSigma_{\X}^{-1} (\t - \Vmu_{\X}) \\
    \hat\VSigma_{pr} &= G_{\Z} - G_{\Z,\X} \hat\VSigma_{\X}^{-1} G_{\Z,\X}^T
    \end{myalign}
    \item With zero mean assumption:
    \begin{myequation}
    \hat\Vmu_{pr} = G_{\Z,\X} \hat\VSigma_{\X}^{-1} \t
    \end{myequation}
    \item For a single test point $\z$:
    \begin{myalign}
    \mu_{pr} &= m(\z) + G_{\z,\X} \hat\VSigma_{\X}^{-1} (\t - \Vmu_{\X}) \\
    \sigma^2_{pr} &= \kappa(\z,\z) - G_{\z,\X} \hat\VSigma_{\X}^{-1} G_{\z,\X}^T
    \end{myalign}
    \item With noise, predictive mean no longer interpolates exactly, and predictive variance at training points is $\sigma_f^2 > 0$.
\end{itemize}

\subsection*{Estimating Kernel Parameters}

\begin{itemize}
    \item Predictive performance depends on suitability of chosen kernel and its hyperparameters.
    \item For RBF kernel:
    \begin{myequation}
    \kappa(\x_i,\x_j)=\sigma^2_f e^{-\frac{1}{2}(\x_i-\x_j)^T\mathbf{M}(\x_i-\x_j)}+\sigma_y^2\delta_{ij}
    \end{myequation}
    \item Simplest $\mathbf{M} = l^{-2}\I$, giving hyperparameters $\sigma_f, \sigma_y, l$.
    \item Varying these parameters significantly affects predictions:
    \begin{itemize}
        \item $l$ (length scale): controls smoothness (larger $l$ $\Rightarrow$ smoother).
        \item $\sigma_f$ (signal variance): controls amplitude of function variations.
        \item $\sigma_y$ (noise variance): controls amount of assumed noise.
    \end{itemize}
    \item Hyperparameters can be estimated by maximizing the log marginal likelihood:
    \begin{myequation}
    \log p(\t|\X) = -\frac{1}{2}\t^T \hat\VSigma_{\X}^{-1} \t - \frac{1}{2}\log|\hat\VSigma_{\X}| - \frac{n}{2}\log 2\pi
    \end{myequation}
    \item This provides an automatic trade-off between data fit and model complexity.
\end{itemize}

\subsection*{Recap of Gaussian Process Properties}

\begin{itemize}
    \item GP provides a nonparametric Bayesian approach to regression.
    \item Directly defines a distribution over functions via mean and covariance functions.
    \item Predictions come with natural uncertainty estimates (predictive variance).
    \item Kernel choice encodes assumptions about function smoothness and structure.
    \item Hyperparameters can be learned from data via marginal likelihood optimization.
    \item GP regression is equivalent to Bayesian linear regression with an infinite number of basis functions.
    \item The equivalent kernel perspective shows GP predictions are weighted averages of training targets, with weights determined by kernel and training data.
    \item GP naturally handles noise and provides interpolation in the noise-free case, regression with uncertainty in the noisy case.
\end{itemize}

